\section{Aftermath: Disengagement and the Road to Peace}

\subsection{Shuttle Diplomacy, Sinai I and Sinai II}
The cease-fire of October 1973 left a tangled military situation crying out for negotiated stabilization. Henry Kissinger seized the moment, embarking on a celebrated campaign of shuttle diplomacy that carried him repeatedly between Middle Eastern capitals. He shuttled among Jerusalem, Cairo, and Damascus, mediating face to face and translating the urgent realities of the battlefield into incremental diplomatic steps. This personal, intensive method allowed him to broker arrangements that the parties could not have reached through direct confrontation, establishing the United States as the central mediator in the region.

The first fruit was the Sinai I disengagement agreement of January 1974, which separated Egyptian and Israeli forces along the canal. Israeli troops withdrew from the west bank and pulled back on the east bank, while a United Nations buffer zone stabilized the front. Crucially, the agreement relieved the encircled Third Army and reduced the immediate danger of renewed fighting. Sinai I demonstrated that the military stalemate could be converted into a managed political process, and it began to rebuild a measure of trust between former combatants.

A parallel disengagement on the Syrian front followed later in 1974, separating forces on the Golan Heights and establishing a monitored zone. Though achieved with greater difficulty, this agreement confirmed that the disengagement model could be extended beyond the Egyptian front. Together these accords reflected Kissinger's step-by-step philosophy, which favored modest, achievable agreements over comprehensive settlements that risked collapse. By breaking an intractable conflict into manageable increments, the approach generated momentum and accustomed the parties to the practice of negotiation.

The Sinai II agreement of September 1975 deepened the Egyptian-Israeli disengagement, providing for further Israeli withdrawal in the Sinai and the establishment of early-warning stations staffed by American personnel. The accord included commitments that bound the United States more closely to the process and signaled Egypt's continued movement away from the Soviet orbit. Sinai II represented a significant advance, demonstrating that incremental diplomacy could produce successive agreements that progressively reduced tensions and laid groundwork for a more comprehensive settlement in the future.

A brief multilateral conference convened at Geneva in December 1973 provided formal cover for the negotiations, even as the real work proceeded through Kissinger's bilateral shuttles. The combination of public framework and private mediation proved effective in advancing the disengagement agreements. By 1975 the pattern was clear: patient, American-brokered diplomacy had transformed a dangerous war into a sequence of stabilizing accords. These step-by-step achievements built the confidence and precedents that would prove essential when the search for a fuller peace gathered force later in the decade. \cite{kissinger1982} \cite{siniver2013} \cite{rabinovich2004}

\subsection{From Sadat's Jerusalem Visit to Camp David and 1979}
The disengagement agreements created a foundation upon which a more ambitious peace could be built, but the decisive breakthrough came from Sadat's audacious personal initiative. In November 1977 the Egyptian president stunned the world by traveling to Jerusalem and addressing the Israeli Knesset. The visit shattered a long taboo against direct contact and dramatized Sadat's willingness to pursue peace openly. It transformed the psychological landscape of the conflict, demonstrating that the limited war of 1973 had indeed opened a path toward negotiation from a position of restored dignity.

Sadat's gesture electrified international opinion and placed enormous pressure on Israeli Prime Minister Menachem Begin to respond. The two leaders, longtime adversaries, found themselves engaged in a process whose momentum few had anticipated. The visit confirmed the central argument that the October war had made peace negotiable by allowing Egypt to bargain as an equal rather than a defeated supplicant. Sadat's psychological parity, won on the battlefield and consolidated through diplomacy, made his dramatic outreach both possible and politically survivable at home.

American mediation proved indispensable in converting the dramatic opening into a durable framework. President Jimmy Carter invited Sadat and Begin to the secluded presidential retreat at Camp David in September 1978, where days of intense and difficult negotiation produced the Camp David Accords. The agreements established a framework for peace between Egypt and Israel and addressed, more ambiguously, the question of Palestinian autonomy. Carter's personal involvement, sustaining the talks through repeated crises, echoed the active brokering role that Kissinger had pioneered in the war's immediate aftermath.

The Camp David framework culminated in the Egyptian-Israeli peace treaty signed in March 1979. Under its terms Israel agreed to withdraw fully from the Sinai Peninsula, and Egypt extended formal recognition and normalized relations. The treaty marked the first peace between Israel and an Arab state, removing the most powerful Arab military from the cycle of confrontation. It represented the fulfillment of the strategic logic Sadat had pursued since 1973, exchanging the leverage gained in war for the recovery of territory through negotiated agreement.

The journey from the surprise attack of October 1973 to the treaty of 1979 vindicates the interpretation of the war as a strategic catalyst rather than merely a military episode. Historians debate how direct the causal chain was and what alternatives might have existed, yet the connection is difficult to dismiss. By shattering the post-1967 stalemate and restoring a sense of parity, the war made peace conceivable. Its enduring legacy lies less in the battlefield outcome than in the diplomatic transformation it ultimately set in motion. \cite{siniver2013} \cite{sadat1978} \cite{rabinovich2004}

\begin{figure}[htbp]
\centering
\includegraphics[width=0.88\linewidth,height=0.58\textheight,keepaspectratio]{figures/ch06_web.jpg}
\caption{The Camp David Accords of 1978, with Sadat, Begin, and Carter, grew out of the diplomatic momentum unleashed by the 1973 war and led to the 1979 peace treaty.}
\label{fig:ch_the_camp_david_accords_of_1978_w}
\end{figure}

\bigskip
\begin{otherlanguage}{hebrew}
\subsection*{סיכום הפרק}

המשא ומתן שלאחר המלחמה הוליד את הסכמי ההפרדה סיני א' וסיני ב' באמצעות הדיפלומטיה המעבורתית של קיסינג'ר. תהליך זה הוביל לביקור סאדאת בירושלים, להסכמי קמפ דייוויד ולחוזה השלום משנת {\beginL 1979\endL}.

\end{otherlanguage}

\clearpage
